\subsubsection{\texorpdfstring{$S_3$}{S3}-闭包曲线的 Jacobian 分解与 Prym 极化：可见因子、点数恒等式与 \texorpdfstring{$A_2$}{A2} 根格刚性}\label{par:xi-terminal-zm-s3-closure-jacobian-prym-a2}
沿用结论 \ref{con:xi-terminal-zm-discriminant-ridge-etale-c3-explicit-model} 的属 \(2\) 曲线 \(C\) 与 \(S_3\)-闭包曲线 \(C_6\)，并令 \(\overline{\mathcal R}\) 为三次预解式平面三次曲线（结论 \ref{con:xi-terminal-zm-resolvent-cubic-plane-cubic-geometry}），其方程为 \(\mathscr R(z,y)=0\)。
下面把该 \(S_3\)-闭包在 Jacobian 与 Prym 层面完全分解为两条可见椭圆因子与一条可见极化刚性。

\begin{theorem}[\texorpdfstring{$S_3$}{S3}-等变同源分解与标准因子的椭圆平方化]\label{thm:xi-terminal-zm-s3-closure-jacobian-splitting}
\leanverified{paper\_xi\_terminal\_zm\_s3\_closure\_jacobian\_splitting}
在上述设定下，存在在 \(\QQ\) 上的同源分解
\[
\boxed{\;
J(C_6)\ \sim_{\QQ}\ J(C)\times J(\overline{\mathcal R})^{2}\; }.
\]
等价地，对任意素数 \(\ell\)，存在 \(G_{\QQ}\)-表示同构
\[
\boxed{\;
H^1_{\acute et}(C_{6,\overline{\QQ}},\QQ_\ell)\ \simeq
H^1_{\acute et}(C_{\overline{\QQ}},\QQ_\ell)\ \oplus
H^1_{\acute et}(\overline{\mathcal R}_{\overline{\QQ}},\QQ_\ell)^{\oplus2}\; }.
\]
等价地，对无分歧循环三重覆盖 \(\pi:C_6\to C\) 的 Prym 簇
\[
\operatorname{Prym}(C_6/C):=\ker\!\bigl(\operatorname{Nm}_{\pi}:J(C_6)\to J(C)\bigr)^{0}
\]
成立
\[
\boxed{\;\operatorname{Prym}(C_6/C)\ \sim_{\QQ}\ J(\overline{\mathcal R})^{2}.\;}
\]
\end{theorem}
\begin{proof}
令 \(X:=C_6\) 并记 \(G:=\operatorname{Aut}_{\QQ}(X/\mathbb P^1_y)\cong S_3\)（结论 \ref{con:xi-terminal-zm-discriminant-ridge-etale-c3-explicit-model}）。
群代数 \(\QQ[G]\) 通过图对应作用在 \(J(X)\) 上，从而给出嵌入
\(
\QQ[G]\hookrightarrow \mathrm{End}^0_{\QQ}(J(X))
\)。
令 \(e_{\mathrm{triv}},e_{\mathrm{sgn}},e_{\mathrm{std}}\in\QQ[G]\) 为 \(S_3\) 的三种不可约表示（平凡、符号、标准二维表示）的中心幂等元，并记
\(
A_{\rho}:=e_{\rho}J(X)
\)。
则有
\[
J(X)\ \sim_{\QQ}\ A_{\mathrm{triv}}\times A_{\mathrm{sgn}}\times A_{\mathrm{std}}.
\]
由于 \(X/G\simeq\mathbb P^1\) 亏格为 \(0\)，故 \(A_{\mathrm{triv}}\sim 0\)，从而
\(
J(X)\sim_{\QQ}A_{\mathrm{sgn}}\times A_{\mathrm{std}}
\)。
\par
对 \(V:=H^0(X,\Omega^1)\) 而言，\(X/G\simeq \PP^1\) 亦给出平凡表示重数为 \(0\)。
又由 \(X/A_3\simeq C\) 且 \(g(C)=2\)，得
\[
\dim_{\QQ}V^{A_3}=2.
\]
并且对任一换位 \(\iota\in S_3\)，由 \(X/\langle\iota\rangle\simeq \overline{\mathcal R}\) 与 \(g(\overline{\mathcal R})=1\) 得
\[
\dim_{\QQ}V^{\langle\iota\rangle}=1.
\]
结合 \(S_3\) 的不可约有理表示分解，可唯一确定
\[
V\ \simeq\ 2\cdot\mathrm{sgn}\ \oplus\ \mathrm{std},
\]
该表示分解与 Jacobian 等型分解完全一致。
\par
记 \(H:=A_3\triangleleft S_3\) 为阶 \(3\) 正规子群，则 \(X/H\simeq C\) 且商映射即 \(\pi:X\to C\)。
因为符号表示在 \(H\) 上平凡而标准表示在 \(H\) 上无不变向量，故 \(H\)-平均幂等元
\[
e_H:=\frac{1}{3}\sum_{h\in H}h
\]
在 \(J(X)\) 上的像恰为 \(A_{\mathrm{sgn}}\)。
另一方面，\(e_H\) 作为“取 \(H\)-不变”的投影，与商映射 \(\pi\) 的 pullback \(\pi^\ast:J(C)\to J(X)\) 在同源意义下给出同一阿贝尔子簇；从而
\(
A_{\mathrm{sgn}}\sim_{\QQ}J(C)
\)。
\par
再令 \(\iota:(y,w,z)\mapsto(y,-w,z)\) 为 \(X\) 的阶 \(2\) 自同构。
其不变子域为 \(\QQ(X)^{\langle\iota\rangle}=\QQ(y,z)\)，故 \(X/\langle\iota\rangle\simeq \overline{\mathcal R}\)。
对应平均投影
\(
e_{\langle\iota\rangle}=\frac12(1+\iota)
\)
在符号表示分量上为零（换位元在符号表示上取 \(-1\)），从而 \(e_{\langle\iota\rangle}J(X)\subseteq A_{\mathrm{std}}\)。
并且 \(e_{\langle\iota\rangle}J(X)\) 在同源意义下等同于商映射 \(X\to\overline{\mathcal R}\) 的 pullback 像，因此给出一个椭圆子簇
\[
E:=\pi_{\iota}^{\ast}J(\overline{\mathcal R})\ \subseteq\ A_{\mathrm{std}},
\qquad
E\sim_{\QQ}J(\overline{\mathcal R}).
\]
最后，由标准表示的有理特征域可得
\(
e_{\mathrm{std}}\QQ[S_3]e_{\mathrm{std}}\cong M_2(\QQ)
\)，从而 \(M_2(\QQ)\subseteq \mathrm{End}^0_{\QQ}(A_{\mathrm{std}})\)。
这迫使 \(A_{\mathrm{std}}\) 在 \(\QQ\)-同源意义下为某椭圆曲线的平方；并取上述非零椭圆子簇 \(E\sim_{\QQ}J(\overline{\mathcal R})\) 即得
\(
A_{\mathrm{std}}\sim_{\QQ}J(\overline{\mathcal R})^2
\)。
合并 \(J(X)\sim_{\QQ}A_{\mathrm{sgn}}\times A_{\mathrm{std}}\) 与 \(A_{\mathrm{sgn}}\sim_{\QQ}J(C)\) 即得所述分解。
\end{proof}

\leanverified{paper\_xi\_terminal\_zm\_s3\_root\_recovery\_coordinate\_automorphisms}
\begin{theorem}[三次预解根的显式复原与 \texorpdfstring{$S_3$}{S3} 自同构的坐标化]\label{thm:xi-terminal-zm-s3-root-recovery-coordinate-automorphisms}
\leanverified{paper\_xi\_terminal\_zm\_s3\_root\_recovery\_coordinate\_automorphisms}
记
\[
\mathscr R(z,y):=z^3+(1+2y)z^2-(1+4y+4y^2)z-(1+5y+13y^2+8y^3),
\]
\[
\Delta(y):=\Disc_z\!\bigl(\mathscr R(z,y)\bigr)=-y(y-1)P_{\mathrm{LY}}(y),
\qquad
P_{\mathrm{LY}}(y)=256y^3+411y^2+165y+32.
\]
在函数域 \(\QQ(C_6)\) 中取一根 \(z=z_1\) 与平方根 \(w\) 使 \(w^2=\Delta(y)\)，并记
\[
\mathscr R'_z(z,y):=\frac{\partial\mathscr R}{\partial z}(z,y).
\]
则其余两根可显式写为
\[
\boxed{\;
z_{2,3}
=-\frac{(1+2y)+z_1}{2}
\pm
\frac{w}{2\,\mathscr R'_z(z_1,y)}\; }.
\]
进一步，定义双有理自同构
\[
\sigma:\ (y,z,w)\dashmapsto
\Bigl(
y,\ -\frac{(1+2y)+z}{2}+\frac{w}{2\,\mathscr R'_z(z,y)},\ w
\Bigr),
\]
\[
\iota:\ (y,z,w)\mapsto (y,z,-w).
\]
则在 \(\QQ(C_6)\) 上满足
\[
\sigma^3=\iota^2=(\iota\sigma)^2=1,
\]
并生成一个与 \(S_3\) 同构的自同构群；其在根集 \(\{z_1,z_2,z_3\}\) 上分别对应三循环与换位。
\end{theorem}
\begin{proof}
令 \(f(T):=\mathscr R(T,y)\in\QQ(y)[T]\)，其三根为 \(z_1,z_2,z_3\)。
由 Vieta 关系，
\[
z_2+z_3=-(1+2y)-z_1.
\]
又
\[
f'(z_1)=\prod_{j\neq1}(z_1-z_j)=(z_1-z_2)(z_1-z_3).
\]
判别式平方根可取为
\[
w=\pm (z_1-z_2)(z_1-z_3)(z_2-z_3),
\qquad
w^2=\Disc(f)=\Delta(y).
\]
故
\[
z_2-z_3=\frac{w}{(z_1-z_2)(z_1-z_3)}=\frac{w}{f'(z_1)}=\frac{w}{\mathscr R'_z(z_1,y)}.
\]
与 \(z_2+z_3\) 联立即得根复原公式。

\(\iota\) 由 \(w\mapsto-w\) 给出，固定 \(y,z\) 且交换 \(z_2,z_3\)；\(\sigma\) 将 \(z_1\) 送往上式中“\(+\)”分支并保持 \(w\) 不变，对应三循环。
由置换关系可得 \(\sigma^3=\iota^2=(\iota\sigma)^2=1\)，故生成群同构于 \(S_3\)。
\end{proof}

\begin{proposition}[二次商映射的分支互补性与六点分支除子]\label{prop:xi-terminal-zm-s3-double-cover-branch-complementarity}
\leanverified{paper\_xi\_terminal\_zm\_s3\_double\_cover\_branch\_complementarity}
令
\[
q:C_6\to \overline{\mathcal R},\qquad (y,w,z)\mapsto (y,z),
\]
并记 \(\pi_y^{(3)}:\overline{\mathcal R}\to\PP^1_y\) 为三次投影。
则 \(q\) 的分支恰由六点组成：
\begin{enumerate}
  \item 对每个分支值 \(y_0\in\{0,1\}\cup\{\alpha_1,\alpha_2,\alpha_3\}\)（\(P_{\mathrm{LY}}(\alpha_i)=0\)），取纤维 \((\pi_y^{(3)})^{-1}(y_0)\) 中唯一的未分歧点 \(Q_{y_0}\)；
  \item 取无穷远纤维中的未分歧点 \(Q_\infty=[2:1:0]\)。
\end{enumerate}
等价地，\(q\) 在 \(\pi_y^{(3)}\) 的分歧点上不分歧，而在其互补点上分歧。
\end{proposition}
\begin{proof}
由定理 \ref{thm:xi-s3-closure-double-cover-explicit-branch-divisor}，\(q\) 的约化分歧支撑正是
\[
Q_0+Q_1+Q_\infty+\sum_{i=1}^3Q_{\alpha_i},
\]
其中 \(Q_0,Q_1,Q_{\alpha_i},Q_\infty\) 分别位于 \(y=0,1,\alpha_i,\infty\) 的纤维中。
另一方面，\(\Disc_z(\mathscr R(z,y))=-y(y-1)P_{\mathrm{LY}}(y)\) 表明
\(\pi_y^{(3)}\) 在每个分支值纤维均为 \((2,1)\) 型，故每条纤维恰有一个分歧点与一个未分歧点。
将该 \((2,1)\) 结构与上述显式分支除子对照，即得“分支--互补”刻画。证毕。
\end{proof}

\begin{corollary}[\texorpdfstring{$C_6$}{C6} 的 Hasse--Weil \texorpdfstring{$L$}{L}-函数分解与有限域点数恒等式]\label{cor:xi-terminal-zm-s3-closure-hasse-weil-factorization}
\leanverified{paper\_xi\_terminal\_zm\_s3\_closure\_hasse\_weil\_factorization}
在定理 \ref{thm:xi-terminal-zm-s3-closure-jacobian-splitting} 的设定下，对任意良素数 \(p\nmid 2\cdot 3\cdot 31\cdot 37\) 有局部因子分解
\[
L_p(C_6,T)=L_p(C,T)\cdot L_p(\overline{\mathcal R},T)^2.
\]
因此全局 Hasse--Weil \(L\)-函数在去掉有限个坏素数 Euler 因子后满足
\[
L(C_6,s)=L(C,s)\cdot L(\overline{\mathcal R},s)^2.
\]
并且对任意有限域 \(k=\FF_q\)（特征为良特征）有点数恒等式
\[
\boxed{\;
\#C_6(k)=\#C(k)+2\,\#\overline{\mathcal R}(k)-2(q+1)\; }.
\]
\end{corollary}
\begin{proof}
\(\QQ\)-同源分解给出对几乎所有素数处的 \(\ell\)-进 Tate 模表示分解
\(
T_\ell J(C_6)\otimes\QQ_\ell\cong T_\ell J(C)\otimes\QQ_\ell\oplus (T_\ell J(\overline{\mathcal R})\otimes\QQ_\ell)^{\oplus 2}
\)，从而良素数处 Frobenius 特征多项式相乘即为局部 \(L\)-因子相乘。
点数恒等式由
\(
\#X(\FF_q)=q+1-\mathrm{Tr}(\mathrm{Frob}_q\mid H^1_{\text{ét}}(X_{\overline{k}},\QQ_\ell))
\)
与迹的可加性直接推出。
\end{proof}

\begin{theorem}[Prym 极化的 \texorpdfstring{$A_2$}{A2} 根格刚性与 \texorpdfstring{$W(A_2)\cong S_3$}{W(A2)=S3} 的极化保持作用]\label{thm:xi-terminal-zm-prym-polarization-a2-rigidity}
令 \(\pi:C_6\to C\) 为结论 \ref{con:xi-terminal-zm-discriminant-ridge-etale-c3-explicit-model} 的次数 \(3\) 无分歧循环覆盖，并令
\[
P:=\operatorname{Prym}(C_6/C)=\ker(\operatorname{Nm}_\pi)^{0}\subset J(C_6).
\]
以 \(J(C_6)\) 的典范主极化限制到 \(P\) 得到 Prym 极化 \(\Xi\)。
则：
\begin{enumerate}
  \item \((P,\Xi)\) 为二维极化阿贝尔簇且 \(\deg(\Xi)=9\)，从而极化型为 \((1,3)\)；
  \item 覆叠的自同构群 \(S_3=\operatorname{Aut}_{\QQ}(C_6/\mathbb P^1_y)\) 诱导忠实同态
  \[
  S_3\hookrightarrow \operatorname{Aut}_{\QQ}(P,\Xi),
  \]
  即 \(S_3\) 在 \(P\) 上以保持 Prym 极化的方式作用；
  \item 在定理 \ref{thm:xi-terminal-zm-s3-closure-jacobian-splitting} 的同源识别 \(P\sim_{\QQ}J(\overline{\mathcal R})^2\) 下，\(\Xi\) 的 Rosati 对称型在 \(J(\overline{\mathcal R})^2\) 上（相对于积极化）可取为 \(\mathrm{GL}_2(\ZZ)\)-同余意义下唯一的 \(S_3\)-不变原始正定二次型，其 Gram 矩阵可选为
  \[
  Q_{A_2}=\begin{pmatrix}2&-1\\-1&2\end{pmatrix}.
  \]
  因而 \(S_3\) 在同调格上的作用同构为 Weyl 群 \(W(A_2)\cong S_3\) 的标准表示。
\end{enumerate}
\end{theorem}
\begin{proof}
由 \(\pi\) 为无分歧次数 \(3\) 覆盖，得 \(\dim P=g(C_6)-g(C)=2\)。
并且对 Jacobian 上的主极化识别，\(\operatorname{Nm}_\pi\) 与 \(\pi^\ast\) 互为对偶同态，且满足经典恒等式
\(
\operatorname{Nm}_\pi\circ\pi^\ast=[3]_{J(C)}
\)。
由 \(\deg([3]_{J(C)})=3^{2\dim J(C)}=3^4=81\) 得 \(\deg(\pi^\ast)=9\)；
因此 \(J(C_6)\) 的主极化在 \(P\) 上的限制亦具有度 \(9\)，即极化型为 \((1,3)\)。
\par
\(S_3\) 作为曲线自同构群作用于 \(J(C_6)\) 并保持其典范主极化。
由于 \(\operatorname{Nm}_\pi\) 与覆盖结构等变，\(P=\ker(\operatorname{Nm}_\pi)^0\) 为 \(S_3\)-不变子簇，限制极化 \(\Xi\) 仍被保持，从而得到 (2)。
\par
最后证 (3)。
由定理 \ref{thm:xi-terminal-zm-s3-closure-jacobian-splitting}，\(P\sim_{\QQ}J(\overline{\mathcal R})^2\) 且 \(\mathrm{End}^0_{\QQ}(P)\cong M_2(\QQ)\)。
在固定同源 \(J(\overline{\mathcal R})^2\to P\) 并以积极化识别对偶后，极化 \(\Xi\) 等价于一个对称正定矩阵 \(Q\in \mathrm{Sym}_2(\QQ)\)，且 \(\deg(\Xi)=\det(Q)^2\)。
由 (1) 得 \(\det(Q)=3\)。
另一方面，\(S_3\) 在 \(P\) 上的作用给出一个二维不可约有理表示（标准表示）\(\rho:S_3\to \mathrm{GL}_2(\QQ)\)，极化保持性等价于
\(
\rho(g)^{\mathsf T}Q\,\rho(g)=Q
\)（\(\forall g\in S_3\)）。
标准表示上 \(S_3\)-不变对称双线性型空间一维（Schur 引理），故 \(Q\) 在同余与正有理倍数意义下唯一；
再以 \(\det(Q)=3\) 固定尺度并取原始整系数代表，即得到 \(Q_{A_2}\)。
\end{proof}

\begin{conjecture}[标准因子的图核极化刚性]\label{conj:xi-terminal-zm-standard-factor-graph-kernel-3torsion}
沿用定理 \ref{thm:xi-terminal-zm-s3-closure-jacobian-splitting} 的记号，令
\[
A_{\mathrm{std}}:=e_{\mathrm{std}}J(C_6),\qquad
E:=J(\overline{\mathcal R}).
\]
则存在定义于 \(\QQ\) 的同源
\[
\varphi:E\times E\longrightarrow A_{\mathrm{std}}
\]
使得其核为反对角 \(3\)-挠图
\[
\ker(\varphi)=\{(P,-P):\ P\in E[3]\}\subset E[3]\times E[3].
\]
并且 \(A_{\mathrm{std}}\) 上由 \(J(C_6)\) 主极化诱导的极化，在该同源识别下为 \((1,3)\)-型。
\end{conjecture}

\begin{conjecture}[判别式二次脊曲线的有理点完全刻画]\label{conj:xi-terminal-zm-discriminant-ridge-rational-points}
令
\[
C:\quad w^{2}=-y(y-1)\bigl(256y^{3}+411y^{2}+165y+32\bigr).
\]
则猜想
\[
\boxed{\;
C(\QQ)=\{\infty,\ (0,0),\ (1,0)\}\; },
\]
其中 \(\infty\) 为该奇次数超椭圆模型的唯一无穷远点。
\end{conjecture}

\begin{remark}[有限高度穷举证据]\label{rem:xi-terminal-zm-discriminant-ridge-rational-points-height-audit}
在有理参数表示 \(y=a/b\)（\(\gcd(a,b)=1\)、\(b>0\)）下，对高度界 \(\max\{\abs a,b\}\le 200\) 的枚举可得到如下仿射点集：
\[
\boxed{\ \{(y,w)\in C(\mathbb Q): \max\{|a|,b\}\le 200\}\;=\;\{(y,w)=(0,0),\ (y,w)=(1,0)\}\ }
\]

可复现核验脚本：\texttt{scripts/exp\_fold\_zm\_discriminant\_ridge\_rational\_points\_height200\_audit.py}。
\end{remark}

\endinput
